%%%%%%%%%%%%%%%%%%%%%%%%%%%%%%%%%
%
%       EXERCÍCIO
%
%%%%%%%%%%%%%%%%%%%%%%%%%%%%%%%%%

\ifdefstring{\atividade}{prova}{%
    \renewcommand{\valorquestao}{\ValorQ\ pontos}
}%

\begin{exercicioBanco}[\valorquestao]
Um estudo registrou o tempo de permanência (em dias) de $10$ pacientes internados em uma ala hospitalar:
\[ 5, \; 7, \; 3, \; 12, \; 7, \; 8, \; 15, \; 7, \; 10, \; 6 \]

Com base nesses dados, determine:

\begin{enumerate}[(a)]
    \item A média, a mediana e a moda do tempo de internação.
    \item O primeiro quartil ($Q_1$) e o terceiro quartil ($Q_3$).
    \item Se o valor discrepante $15$ fosse substituído por $25$, qual das três medidas do item (a) sofreria maior alteração? Justifique brevemente.
\end{enumerate}
\end{exercicioBanco}
